% Direct-authority diplomatic French for printed p.135.
\SourcePage{135}

\LayerTarget{tohoku:prop:1.9.2}
\noindent\textsc{Proposition 1.9.2.} --- \textit{Soit $\Sigma$ un schéma de
diagrammes avec relations de commutation (cf. 1.6), et soient $C$ une
catégorie abélienne, $(U_i)_{i\in I}$ une famille de générateurs de $C$. On
suppose que pour toute flèche de $\Sigma$, l'origine et l'extrémité de la
flèche sont distincts, et que dans les relations de commutation ne figurent
pas les morphismes identiques $e_s$ (où $s$ est un sommet)\textsuperscript{2)}.
Alors pour tout $A\in C$, et tout sommet $s$ du schéma, le diagramme
$\varepsilon_s(A)$ dont la valeur est $A$ au sommet $s$ et $0$ en tous les
autres sommets, et dont la valeur sur chaque flèche est $0$, appartient à
$C^\Sigma$. De plus, la famille des $\varepsilon_s(U_i)$ ($s$ et $i$
variables) est un système de générateurs de $C^\Sigma$.}

\LayerTarget{tohoku:prop:1.9.2:proof}
La vérification est immédiate, il suffit de remarquer, pour la dernière
assertion, que les morphismes de $\varepsilon_s(A)$ dans un diagramme $D$
s'identifient aux morphismes de $A$ dans $D(s)$.

\LayerTarget{tohoku:def:cogenerator-family}
Nous laissons au lecteur le soin de développer la notion duale de
\textit{famille de cogénérateurs} d'une catégorie abélienne. On peut montrer
que si $C$ est une catégorie abélienne satisfaisant à l'axiome AB 5)
(cf. 1.5.) alors l'existence d'un générateur implique l'existence d'un
cogénérateur (nous ne nous servirons pas de ce résultat). Ainsi la catégorie
des modules unitaires à gauche sur un an neau $U$ avec unité admet toujours un
cogénérateur : si par exemple $U=\mathbf Z$, on peut prendre pour cogénérateur
le groupe des nombres rationnels mod $\mathbf Z$ (ou le tore
$\mathbf T=\mathbf R/\mathbf Z$).

\LayerTarget{tohoku:section:1.10}
\noindent\textbf{1.10. Objets injectifs et projectifs.}

\LayerTarget{tohoku:def:injective-object}
Rappelons qu'un objet $M$ d'une catégorie abélienne $C$ est dit
\textit{injectif} si le foncteur $A\longmapsto\mathrm{Hom}(A,M)$ (qui de toutes
façons est exact à gauche) est exact, i. e. si pour tout morphisme $u$ dans
$M$ d'un sous-truc $B$ d'un $A\in C$, il existe un morphisme de $A$ dans $M$
qui le prolonge.

\LayerTarget{tohoku:def:injective-effacement}
Un morphisme $A\longrightarrow M$ est dit un \textit{effacement injectif} (de
$A$) si c'est un monomorphisme, et si pour tout monomorphisme
$B\longrightarrow C$ et tout morphisme $B\longrightarrow A$, on peut trouver
un morphisme $C\longrightarrow M$ rendant commutatif le diagramme

\LayerTarget{tohoku:diagram:injective-effacement}
\[
\begin{array}{ccc}
B & \longrightarrow & C \\
\big\downarrow && \big\downarrow \\
A & \longrightarrow & M
\end{array}
\]

Ainsi, pour que l'application identique de $M$ soit un effacement injectif il
faut et il suffit que $M$ soit injectif ; tout monomorphime dans un objet
injectif est un effacement injectif.

\LayerTarget{tohoku:thm:1.10.1:p1}
\noindent\textsc{Théorème 1.10.1.} --- \textit{Si $C$ satisfait à axiome
AB 5) (cf. 1.5.) et admet un générateur (cf. 1.9) alors pour tout $A\in C$, il
existe un monomorphisme de $A$ dans un objet injectif $M$.}

On va même construire un foncteur
$M:A\longrightarrow M(A)$ (non additif en général !)

\LayerTarget{tohoku:footnote:2}
\begin{quote}
\footnotesize
2) La catégorie $C^\Sigma$ des systèmes inductifs dans $C$ construits sur un
ensemble d'indices \textit{ordonné} $I$, rentre dans ce cas. En effet il suffit
alors, dans l'exemple 1.7.h) de considérer les flèches $(i,j)$ avec $i<j$, et
les relations de commutation $(i,j)(j,k)=(i,k)$, où les $e_s$ n'interviennent
plus.
\end{quote}

\TohokuContinuation{136}
